\documentclass[11pt]{article}

\usepackage[a4paper,margin=1in]{geometry}
\usepackage{amsmath,amssymb,amsthm,mathtools}
\usepackage{booktabs,longtable,array}
\usepackage{enumitem}
\usepackage{microtype}
\usepackage[hidelinks]{hyperref}

\newcommand{\R}{\mathbb{R}}
\newcommand{\E}{\mathbb{E}}
\newcommand{\tr}{\operatorname{tr}}
\newcommand{\rank}{\operatorname{rank}}
\newcommand{\range}{\operatorname{range}}
\newcommand{\diag}{\operatorname{diag}}
\newcommand{\pinv}{^{\dagger}}
\newcommand{\norm}[1]{\lVert #1\rVert}
\newcommand{\transpose}{^{\mathsf T}}

\title{\textbf{A Literature Survey and Design Recommendation for\
Robust Direct-Factor Square-Root Sigma-Point Filtering}}
\author{BayesFilter research note}
\date{17 August 2026}

\begin{document}
\maketitle

\begin{abstract}
This note surveys square-root Kalman, unscented, cubature, rank-revealing, and
matrix-differentiation methods relevant to a sigma-point filter that must remain
usable for ill-conditioned or structurally singular covariance matrices and must
not form a covariance followed by an eigendecomposition or singular value
decomposition at every step.  The conclusion is deliberately narrower than the
phrase ``SVD is more stable.''  Classical square-root filtering does exactly the
right arithmetic operation: it factors weighted residual stacks directly and
uses triangular solves.  The most useful construction for the present problem is
a rectangular loading-factor filter with a fixed-rank, fixed-pivot QR branch for
value-and-score computation.  A rank-revealing QR or SVD is retained as a
diagnostic and as a value-only support/rank-selection operation.  For a
singular innovation, the likelihood must be a density on its affine support,
with a pseudodeterminant and a range-consistency test; an ordinary full-space
Gaussian likelihood is not defined.  A block QR of a joint residual stack gives
the conditional filtered factor directly and removes the need for covariance
subtraction or sequential downdates.  This construction is source-faithful to
the square-root filtering literature where stated, and is identified as an
extension where the literature does not supply the exact singular,
score-bearing sigma-point algorithm.
\end{abstract}

\tableofcontents

\section{Question, scope, and short answer}

The question is whether the current direct-factor square-root unscented Kalman
filter (SR--UKF) idea is anchored in the literature, and which method should be
chosen when covariance matrices are badly conditioned or singular.  Four cases
must be kept separate:
\begin{enumerate}[label=(\roman*)]
  \item a positive-definite covariance with a large condition number;
  \item a positive-semidefinite covariance with a known structural rank defect;
  \item a singular innovation covariance, which changes the probability measure
        and therefore the likelihood formula; and
  \item a parameter path on which rank, a pivot pattern, or a spectral
        multiplicity changes, which is a nonsmooth algorithmic branch for score
        computation.
\end{enumerate}

The answer is not to replace every QR operation by an SVD.  The literature gives
the following division of labour.

\begin{itemize}
  \item Classical square-root filters propagate a factor and obtain it from
        orthogonal triangularization of a residual stack, avoiding the explicit
        formation and square-rooting of a covariance matrix
        \cite{kaminski1971, bierman1977, maybeck1979, vandermerwewan2001,
        arasaratnamhaykin2008}.
  \item An SVD is an excellent rank and conditioning diagnostic, and a robust
        value-only factorization when the active rank is allowed to change.  It
        is not a globally smooth coordinate system: repeated singular values
        make singular vectors nonunique, and rank changes make the truncated
        factor discontinuous as a map of the input matrix
        \cite{stewartsun1990, golubvanloan2013, ionescu2015}.
  \item Fixed-order or fixed-pivot QR has a locally differentiable triangular
        gauge whenever the selected stack has fixed rank and nonzero pivots.
        Pivot changes and zero pivots are branch events, not details that can be
        differentiated through silently \cite{magnusneudecker1999,
        golubvanloan2013, chan1987, gu1996}.
  \item The proposed robust design therefore carries a rectangular factor
        directly, uses QR for the admitted differentiable branch, and calls
        RRQR/SVD only to detect or report a rank/support change.  A singular
        innovation is handled in its support coordinates, not by adding an
        unreported nugget.
\end{itemize}

The block-QR conditional update derived in Section~\ref{sec:conditional} is a
particularly useful refinement.  It is a direct factor update, not a covariance
subtraction disguised as a factorization.

\section{Filtering problem and numerical requirements}

Consider the additive-noise state-space model
\begin{equation}
 x_t=f_t(x_{t-1},q_t;\theta),\qquad q_t\sim N(0,Q_t),
 \qquad y_t=h_t(x_t;\theta)+r_t,\qquad r_t\sim N(0,R_t).
 \label{eq:model}
\end{equation}
The exact nonlinear filter is generally intractable.  The UKF of Julier and
Uhlmann~\cite{julieruhlmann1997} and the broader sigma-point treatment of van
der Merwe~\cite{vandermerwe2004} replace Gaussian moment integrals by a finite
weighted point rule.  The rule is an approximation to moments; square-root
bookkeeping does not make the resulting likelihood exact.

The numerical requirements here are stronger than merely returning a finite
number.  A candidate implementation should:
\begin{enumerate}
  \item carry a factor $G_t\in\R^{n\times r}$ with
        $P_t=G_tG_t\transpose$ and allow $r<n$;
  \item form factors from residual/loadings directly, never by first forming
        $P_t$ and calling $\operatorname{chol}(P_t)$, $\operatorname{eigh}(P_t)$,
        or $\operatorname{svd}(P_t)$ at each time step;
  \item use solves rather than explicit inverses;
  \item distinguish an exact positive-definite branch, a rank-truncated
        positive-semidefinite branch, and a failed/unsupported branch; and
  \item expose enough pivots, singular values, support residuals, and
        reconstruction residuals to audit a long run.
\end{enumerate}

The first item is important.  A square factor is not required by the probability
law: if $G$ is rectangular, then $GG\transpose$ is still a valid covariance,
and a factor is nonunique under right orthogonal transformations.  The filter
must therefore specify a factor gauge (triangular, fixed QR, or a fixed support
chart) as part of its algorithm identity.

\section{What the classical literature actually establishes}

\subsection{Linear square-root filtering}

The early square-root filtering literature was motivated by two familiar
failure modes of covariance-form recursions: roundoff destroys symmetry and
positive semidefiniteness, and subtracting two nearly equal positive matrices
loses significant digits.  Kaminski, Bryson, and Schmidt's survey
\cite{kaminski1971} and Bierman's factorization treatment
\cite{bierman1977} establish the central pattern: maintain a factor, use
orthogonal transformations or rank-one updates, and perform measurement solves
in factor coordinates.  Maybeck~\cite{maybeck1979} gives the standard
state-space estimation context.  These are numerical representations of the
same Kalman algebra, not different statistical models.

For a linear prediction $x^-=Fx^++w$ with factors
$P^+=G^+G^{+\transpose}$ and $Q=G_QG_Q\transpose$, the direct stack is
\begin{equation}
 A^-=[\,FG^+\ \ G_Q\,].
 \label{eq:linear-pred-stack}
\end{equation}
If $A^{-\transpose}=QR$, with a positive diagonal convention, then
$G^-=R\transpose$ and
$G^-G^{-\transpose}=A^-A^{-\transpose}$.  No covariance is constructed.
The measurement stack is analogous.  This is exactly the operation that the
present problem needs.

\subsection{Sigma-point filters and the SR--UKF}

Julier and Uhlmann~\cite{julieruhlmann1997,julieruhlmann2004} introduced the
unscented transformation and its filtering use.  Van der Merwe and Wan's
ICASSP paper~\cite{vandermerwewan2001} explicitly describes the square-root
UKF: the factor is propagated using QR decomposition, Cholesky updates, and
pivot-based least-squares solves.  Their thesis gives the wider sigma-point
family and implementation context~\cite{vandermerwe2004}.  Thus a direct
factor SR--UKF is not an invention without precedent.

There is, however, a qualification that matters for robust rank-deficient work.
Many scaled UKF parameterizations assign a negative covariance weight to the
central point.  A positive residual stack then represents only the positive
part, and the negative part requires a downdate.  A downdate can fail when the
result is not positive definite.  This is why a ``pseudo square-root UKF'' is
not automatically robust for arbitrary UKF weights.  The present DZ5 rule with
$\alpha=1$, $\beta=2$, $\kappa=0$ has
\begin{equation}
 w_0^{(m)}=0,\qquad w_0^{(c)}=2,\qquad
 w_i^{(m)}=w_i^{(c)}=\frac{1}{2d},\quad i=1,\ldots,2d,
 \label{eq:dz5-weights}
\end{equation}
so all covariance weights are nonnegative.  That special choice permits a
pure residual-stack QR implementation.  It must not be silently generalized to
negative-weight UKF rules.

\subsection{Square-root cubature filtering}

Arasaratnam and Haykin's CKF paper~\cite{arasaratnamhaykin2009} derives a
third-degree cubature filter and a square-root CKF.  Their square-root section
explicitly identifies matrix square-rooting, inversion, squared-form roundoff,
and covariance subtraction as sensitive operations; it replaces them by
triangularization, least-squares solves, and factor updates.  The square-root
CKF is therefore a particularly clear precedent for direct residual-stack
factorization.  Their 2008 square-root CKF paper~\cite{arasaratnamhaykin2008}
also gives the compact algorithmic form.

The CKF and UKF use different point rules.  The numerical lesson transfers, but
one should not claim that the CKF derivation proves a singular SR--UKF theorem.
It proves that the direct-factor architecture is an established square-root
filtering pattern.

\subsection{Adjacent robust-filter literature}

There is a second meaning of ``robust UKF'' in the filtering literature.  Yang,
Shi, and Chen~\cite{yang2019robust} use generalized maximum-likelihood
M-estimation to downweight functional-model and measurement outliers.  Tseng,
Lin, and Jwo~\cite{tseng2017robust} similarly combine Huber M-estimation with
the CKF for contaminated GPS observations.  These are useful references for an
outlier-robust observation model, but they do not solve the numerical problem
studied here: neither paper establishes a direct rectangular factor recursion
through a rank-deficient covariance, a support-density likelihood, or a smooth
derivative through a rank transition.  Their residual reweighting can be
composed with the proposed filter only as a declared change to the statistical
model.  It must not be presented as evidence that SVD resolves singular
square-root filtering.

Koch and Yang's rank-deficient-observation Kalman filter~\cite{kochyang1998}
is a closer statistical neighbour.  It shows that a rank-deficient observation
model can be treated as a real estimation problem rather than automatically
repaired by an arbitrary positive diagonal.  It does not, however, supply the
nonlinear sigma-point, direct residual-stack, or differentiable support-chart
construction proposed here.

For BayesFilter, the locally verified primary sources are the Julier--Uhlmann
paper in \path{.local_sources/highdim_nonlinear_filtering/}, the van der Merwe
thesis PDF under \path{docs/}, and the Arasaratnam--Haykin CKF PDF in the same
local source directory.  The repository's existing derivations in
\path{docs/fable-rewrite/monograph/chapters/ch17_square_root_sigma_point.tex},
\path{ch18_svd_sigma_point.tex}, and \path{ch12_factor_derivatives.tex} were
used as implementation cross-checks, not as substitutes for the external
literature.

\begin{table}[ht]
\centering
\caption{Literature status of the main ingredients.}
\label{tab:literature-status}
\begin{tabular}{p{0.28\linewidth}p{0.30\linewidth}p{0.31\linewidth}}
\toprule
Ingredient & What is established & Boundary for this project \\
\midrule
Sigma-point moments & UKF/UT and sigma-point filtering
  \cite{julieruhlmann1997,julieruhlmann2004,vandermerwe2004} & Moment approximation,
  not exact nonlinear likelihood \\
Square-root recursion & QR/orthogonal triangularization and factor solves
  \cite{kaminski1971,bierman1977,vandermerwewan2001,arasaratnamhaykin2009} & Usually
  assumes a regular positive factor or treats downdates separately \\
Rank revelation & RRQR and SVD identify numerical rank
  \cite{chan1987,gu1996,golubvanloan2013} & Rank selection is a discrete branch \\
Matrix derivatives & Fixed-rank QR and matrix perturbation calculus
  \cite{magnusneudecker1999,stewartsun1990,ionescu2015} & No globally smooth factor gauge at
  rank changes or repeated singular values \\
Degenerate Gaussian likelihood & Support/range and generalized-inverse
  formulations are standard measure-theoretic linear algebra
  \cite{bogachev1998,benisraelgreville2003} & Must be implemented explicitly; an ordinary
  ambient Gaussian density is invalid \\
\bottomrule
\end{tabular}
\end{table}

\section{Direct residual-stack factors}

Let $X=[x_1-\bar x,\ldots,x_m-\bar x]$ be a centered state residual stack and
let $w_j\geq0$ be covariance weights.  Include additive noise loadings as
columns of $G_Q$.  Define
\begin{equation}
 A_x=[\,\sqrt{w_1}(x_1-\bar x)\ \cdots\ \sqrt{w_m}(x_m-\bar x)\ \ G_Q\,].
 \label{eq:state-stack}
\end{equation}
Then the covariance represented by the point rule is exactly
$P_x=A_xA_x\transpose$.  Apply a thin QR to $A_x\transpose$:
\begin{equation}
 A_x\transpose=Q_xR_x,\qquad G_x=R_x\transpose,
 \qquad G_xG_x\transpose=A_xA_x\transpose.
 \label{eq:prediction-qr}
\end{equation}
The factor is generated from the stack, not from its Gram matrix.  The same
construction applies to the innovation stack
\begin{equation}
 A_y=[\,\sqrt{w_j}(y_j-\bar y)\ \cdots\ \ G_R\,].
 \label{eq:innovation-stack}
\end{equation}
Parameter-dependent observation noise must be differentiated as a loading
column, not reconstructed from a differentiated covariance.

For a rectangular prior factor $G_a\in\R^{n_a\times r_a}$, use latent offsets
$\xi_j\in\R^{r_a}$ and points
\begin{equation}
 a_j=\mu_a+G_a\xi_j.
 \label{eq:rectangular-sigma-points}
\end{equation}
This is the natural reduced-rank generalization of sigma points.  It preserves
the represented law and never asks for a square root of $G_aG_a\transpose$.
The point rule must state whether its offsets are defined in factor coordinates
or after an additional orthogonal rotation; the two choices can produce
different nonlinear moment approximations even though they have the same
covariance.

\section{The conditional block-QR update}
\label{sec:conditional}

The usual covariance-form update is
\begin{equation}
 P_f=P^- - P_{xy}S^{-1}P_{yx},\qquad K=P_{xy}S^{-1}.
 \label{eq:cov-update}
\end{equation}
Computing the subtraction and refactorizing it is precisely what a direct
square-root implementation should avoid.  There is a direct alternative.

Build one joint residual/loading stack, with observation rows first.  If
$\widetilde Y$ and $\widetilde X$ are the weighted point residual columns,
the observation-noise loading is padded with a zero state block:
\begin{equation}
 A=\left[
 \begin{array}{c|c}
 \widetilde Y & G_R\\
 \widetilde X & 0
 \end{array}\right]
 =\begin{bmatrix}Y\\X\end{bmatrix},
 \qquad A\transpose=Q
 \begin{bmatrix}R_{yy}&R_{yx}\\0&R_{xx}\end{bmatrix},
 \label{eq:block-qr}
\end{equation}
where the block QR is taken with a fixed row/column ordering and a positive
diagonal convention.  Since $A A\transpose=R\transpose R$,
\begin{align}
 S&=R_{yy}\transpose R_{yy},\\
 P_{xy}&=R_{yx}\transpose R_{yy},\\
 P^-&=R_{yx}\transpose R_{yx}+R_{xx}\transpose R_{xx}.
 \label{eq:block-covariance-identities}
\end{align}
Consequently,
\begin{equation}
 G_f=R_{xx}\transpose,
 \qquad G_fG_f\transpose=P^- - P_{xy}S^{-1}P_{yx},
 \label{eq:conditional-factor}
\end{equation}
and, without forming an inverse,
\begin{equation}
 K=R_{yx}\transpose R_{yy}^{-\transpose}.
 \label{eq:block-gain}
\end{equation}
The lower-right block is the conditional residual factor.  This is a Schur
complement computed by orthogonal elimination, not a covariance subtraction
followed by a new square root.  It is the recommended filtered-factor update
for nonnegative covariance weights and a full-rank innovation block.

The process-noise loading $G_Q$ appears in the prediction stack
~\eqref{eq:state-stack}; it is not appended a second time here.  The state
residual columns $\widetilde X$ already represent the predicted state law.

The equivalence follows from the block Gram identity.  It is also a useful unit
test: compare the reconstructed covariance only at the test boundary, while
the runtime carries $R_{xx}\transpose$ and never constructs $P_f$.

For negative sigma-point weights, put positive columns and negative columns in
separate stacks.  The difference of their Gram matrices is not itself a single
ordinary residual stack.  Classical SR--UKF implementations use ordered
Cholesky updates/downdates for this case~\cite{vandermerwewan2001}.  The block
QR formula is therefore the primary route for the nonnegative DZ5 rule; a
general signed-weight route needs a separately audited hyperbolic or downdate
factorization and must fail closed when the result leaves the PSD cone.

\section{QR, Cholesky, RRQR, and SVD: a precise comparison}

\subsection{Full-rank ill-conditioning}

For a known full-rank stack, QR applies orthogonal transformations and avoids
forming the normal equations $A A\transpose$.  This is the relevant stability
advantage.  A Cholesky factorization of $A A\transpose$ squares the condition
number before factoring; this is exactly the operation the square-root
literature seeks to avoid.  A QR factor is triangular and therefore convenient
for solves and fixed-branch differentiation.

SVD is also backward stable and gives singular values explicitly.  It is often
the best way to measure a condition number and to decide whether a pivot is
numerically negligible.  But if the only input to SVD is the already formed
covariance $P=A A\transpose$, the algorithm has lost the main square-root
benefit: the Gram matrix has already squared the conditioning and may already
have lost small directions.  An SVD of the residual stack $A$ is a legitimate
direct square-root operation; an SVD of $P$ is not the desired architecture.

\subsection{Rank deficiency and rank revelation}

Ordinary unpivoted QR assumes the selected columns are independent.  If a
stack is rank deficient, a triangular diagonal is zero and the factor gauge is
not unique.  Column-pivoted QR and strong RRQR~\cite{chan1987,gu1996} expose a
numerical rank while retaining orthogonal transformations.  A direct SVD of the
stack $A$ gives
\begin{equation}
 A=U_r\Sigma_rV_r\transpose,
 \qquad G=U_r\Sigma_r,
 \qquad GG\transpose=A A\transpose,
 \label{eq:svd-stack}
\end{equation}
on a fixed retained rank $r$.  This is mathematically valid and does not form
$A A\transpose$.

The price is a branch.  The retained rank depends on a tolerance, and the
singular vectors are nonunique under repeated singular values.  The factor
$U_r\Sigma_r$ can rotate discontinuously even when the covariance changes
smoothly.  A value-only filter can accept this and record the rank and spectrum.
An analytical score cannot claim one smooth derivative across the event.

The recommended implementation consequently has two explicit modes:
\begin{description}
  \item[Fixed-rank score mode.] Use a fixed-order QR or a fixed pivot pattern
    established before the differentiated run.  Require positive pivots above
    a scale-aware threshold and fail closed otherwise.
  \item[Rank-revealing value mode.] Apply RRQR or SVD directly to each residual
    stack, retain a declared rank, and emit the truncation tolerance, spectrum,
    support basis, and discarded-tail norm.  This mode is not a smooth score
    route unless a separate fixed-support derivative contract is proved.
\end{description}

\subsection{A fixed-pivot rectangular QR chart}

The fixed-rank score mode can be made constructive without an SVD at every
step.  Let $B=A\transpose\in\R^{k\times n}$ have exact rank $r<n$.  A
rank-revealing preflight chooses a fixed state-coordinate permutation $\Pi$ so
that
\begin{equation}
 B\Pi=[\,B_1\ \ B_2\,],\qquad B_1\in\R^{k\times r}
 \label{eq:pivot-partition}
\end{equation}
has full column rank.  On this fixed chart, compute only the ordinary thin QR
\begin{equation}
 B_1=Q_1R_{11},\qquad R_{12}=Q_1\transpose B_2,
 \qquad \widehat R=[\,R_{11}\ \ R_{12}\,].
 \label{eq:rectangular-qr-chart}
\end{equation}
If $B_2\in\range(B_1)$, then $B\Pi=Q_1\widehat R$ exactly.  The rectangular
covariance factor is
\begin{equation}
 G=\Pi\widehat R\transpose\in\R^{n\times r},
 \qquad GG\transpose=A A\transpose.
 \label{eq:rectangular-qr-factor}
\end{equation}
No zero diagonal is passed to a square QR derivative and no singular vectors
are selected.  The only differentiated QR is the full-column-rank matrix
$B_1$.  In particular,
\begin{equation}
 \dot R_{12}=\dot Q_1\transpose B_2+Q_1\transpose\dot B_2,
 \qquad
 \dot G=\Pi[\,\dot R_{11}\ \ \dot R_{12}\,]\transpose,
 \label{eq:rectangular-qr-derivative}
\end{equation}
where $(\dot Q_1,\dot R_{11})$ follow the ordinary QR formulas in
Section~\ref{sec:qr-derivative}.

In finite precision, record the chart residual
\begin{equation}
 E_2=(I-Q_1Q_1\transpose)B_2.
 \label{eq:chart-residual}
\end{equation}
If $E_2$ is nonzero above tolerance, the factor reconstructs the projected
stack, not the requested stack.  Score mode must fail closed; value mode may
repivot or invoke direct-stack SVD and report a rank change.  This construction
is a direct application of rank-revealing QR ideas~\cite{chan1987,gu1996}; the
particular factor/derivative contract in
equations~\eqref{eq:pivot-partition}--\eqref{eq:chart-residual} is the proposed
BayesFilter specialization.

\begin{table}[ht]
\centering
\caption{Choosing a factorization for the proposed filter.}
\begin{tabular}{p{0.21\linewidth}p{0.26\linewidth}p{0.25\linewidth}p{0.20\linewidth}}
\toprule
Method & Strength & Failure/branch event & Recommended role \\
\midrule
Cholesky of $P$ & Simple triangular solves & Forms Gram matrix; fails at PSD boundary &
  Strict-SPD reference only \\
Unpivoted QR of $A\transpose$ & Direct stack factor; local smooth triangular gauge &
  Zero pivot or rank loss & Fixed-rank score branch \\
Pivoted/RRQR of $A\transpose$ & Rank estimate without Gram matrix & Pivot pattern/rank changes &
  Value route and diagnostics \\
SVD of $A$ & Best spectrum/support diagnostic; rectangular factor & Repeated values,
  rank threshold, vector gauge & Value route and independent comparator \\
SVD of $P$ & Spectral solve for an existing covariance & Has already formed $P$ and
  squares conditioning & Historical/diagnostic covariance backend \\
\bottomrule
\end{tabular}
\end{table}

\section{Singular innovations are a measure-theoretic change}

Suppose the innovation covariance has rank $r_y<n_y$ and a thin support factor
$G_y=U_y\Sigma_y$, where $U_y\transpose U_y=I_{r_y}$ and $\Sigma_y$ is positive
diagonal.  The Gaussian law is concentrated on
$\bar y+\range(U_y)$.  Let $e=y-\bar y$.  The correct support checks are
\begin{equation}
 e_\perp=(I-U_yU_y\transpose)e,
 \qquad \norm{e_\perp}\leq\tau_{\rm support}\,\max(1,\norm e).
 \label{eq:support-check}
\end{equation}
If the check fails, the observation has zero probability under the degenerate
law.  If it passes, define $z=\Sigma_y^{-1}U_y\transpose e$ and use the
support density
\begin{equation}
 \ell=-\frac12\left[r_y\log(2\pi)+2\sum_{j=1}^{r_y}\log\Sigma_{y,jj}
 +z\transpose z\right].
 \label{eq:singular-likelihood}
\end{equation}
Equivalently, this uses the pseudodeterminant and Moore--Penrose inverse on the
support.  It is not the $n_y$-dimensional Lebesgue density with a zero
eigenvalue.  Degenerate Gaussian measures and generalized inverses are treated
in this form in the probability and linear-algebra literature
\cite{bogachev1998,benisraelgreville2003}.

This distinction is essential for robust filtering.  Adding $\varepsilon I$ can
be a declared model regularization, but it changes the likelihood and its
derivative.  It must not be presented as a numerically invisible fix for
singularity.  A support-aware route should record rank, smallest retained
singular value, support residual, pseudodeterminant, and whether the observation
was rejected.

\section{Analytical derivatives and the equal-singular-value problem}

\subsection{Fixed-branch QR differentiation}
\label{sec:qr-derivative}

Let $B(\theta)=Q(\theta)R(\theta)$ be a thin QR with full column rank and
$R_{jj}>0$.  For a perturbation $\dot B$, set
\begin{equation}
 E=Q\transpose\dot B R^{-1},
 \qquad \Omega=\operatorname{sl}(E)-\operatorname{sl}(E)\transpose,
 \label{eq:qr-derivative-def}
\end{equation}
where $\operatorname{sl}$ retains the strict lower triangle.  Then
\begin{align}
 \dot R&=(E-\Omega)R,\\
 \dot Q&=Q\Omega+(I-QQ\transpose)\dot B R^{-1}.
 \label{eq:qr-derivative}
\end{align}
The positive diagonal fixes the sign gauge locally.  The equations follow by
differentiating $B=QR$ and $Q\transpose Q=I$; they are standard matrix
calculus, not a heuristic finite difference~\cite{magnusneudecker1999,
stewartsun1990}.  The derivatives reconstruct
\begin{equation}
 \dot B=\dot Q R+Q\dot R,
 \qquad
 \dot G G\transpose+G\dot G\transpose=\dot B\transpose B+B\transpose\dot B
 \label{eq:qr-reconstruction}
\end{equation}
for $G=R\transpose$.

The assumptions matter.  If a pivot reaches zero, or if a pivoting policy
selects a different column, equations~\eqref{eq:qr-derivative} describe a
different local chart.  A score-bearing implementation should report a branch
failure rather than differentiating through a hidden pivot or sign change.

\subsection{Why SVD derivatives are not a universal escape}

For $B=U\Sigma V\transpose$, simple, separated singular values permit local
perturbation formulas.  At a repeated singular value, however, the associated
singular subspace is defined but its individual vectors are not.  Any rotation
inside that subspace is an equally valid SVD.  At a rank threshold, the retained
subspace itself changes.  Thus a derivative of $U\Sigma$ requires a selected
gauge and a spectral-gap assumption.  This is the ``equal eigenvalue'' problem
in SVD clothing.  Matrix backpropagation surveys and perturbation theory make
the same point~\cite{stewartsun1990,ionescu2015}.

The physical covariance $BB\transpose$ may have a perfectly meaningful
directional derivative even when a particular SVD factor does not.  That does
not rescue a score route whose sigma points depend on the chosen factor: a
rotated factor produces a different nonlinear point cloud.  The clean policy is
therefore:
\begin{itemize}
  \item do not use an SVD factor derivative across a repeated-value or rank
        transition;
  \item use fixed-rank QR for the differentiable branch; and
  \item use SVD/RRQR to detect a branch event and to provide an independent
        value-level rank diagnostic.
\end{itemize}

\section{Literature-grounded proposed algorithm}
\label{sec:algorithm}

The following algorithm is the recommended target for BayesFilter.  It is
source-faithful to the classical square-root and SR--UKF operations, while the
rectangular singular-support and block-conditional pieces are explicitly
identified as a construction from those ingredients.

\subsection{State representation and points}

Carry $(\mu_t,G_t)$ with $G_t\in\R^{n_x\times r_t}$ and
$P_t=G_tG_t\transpose$.  For fixed standardized offsets $\xi_j\in\R^{r_t}$,
place points by equation~\eqref{eq:rectangular-sigma-points}.  For additive
process noise, either augment the latent factor with $G_Q$ or append its loading
columns to the prediction residual stack, exactly once.  The choice must be
fixed by the route identity; appending noise twice is an implementation error.

Propagate points through $f_t$ and $h_t$.  Compute means and centered residual
stacks in TensorFlow, preserving batch and derivative axes.  The value path uses
equations~\eqref{eq:state-stack}--\eqref{eq:prediction-qr}; it never creates
$P^-$.  In score mode, use the fixed QR derivative equations and check the
factor reconstruction residual.

\subsection{Prediction and innovation}

Form the observation stack with observation-noise loadings.  For a full-rank
innovation, use the block arrangement in equation~\eqref{eq:block-qr}.  The
lower-right block supplies the filtered factor directly, and the off-diagonal
block supplies the gain.  The log likelihood uses triangular solves and the
diagonal of $R_{yy}$.

For a rank-deficient innovation, run a rank-revealing factorization on the
observation stack itself.  Freeze the selected rank and support chart for a
value-and-score window only after checking a spectral/pivot gap.  If the chart
is not fixed, use the value-only support likelihood or fail closed for the
analytical score.  A possible implementation is:
\begin{enumerate}
  \item compute a direct SVD $A_y=U_y\Sigma_yV_y\transpose$ (an RRQR of the
        stack may be used as a cheaper rank diagnostic) to obtain the support
        basis $U_y$, retained singular values, and $r_y$;
  \item support consistency test~\eqref{eq:support-check};
  \item support-coordinate solve and pseudodeterminant likelihood
        ~\eqref{eq:singular-likelihood};
  \item block elimination on the support coordinates to obtain the conditional
        state factor; and
  \item telemetry for rank, gap, discarded tail, and support residual.
\end{enumerate}

\subsection{Branch policy}

The branch identity must bind: point offsets and weights, factor orientation,
QR pivot policy, positive-diagonal convention, retained rank, support basis,
regularization policy, and likelihood measure.  A pivot sign correction,
rank change, active floor, support rejection, or downdate failure is recorded
as a branch event.  It is not silently repaired by adding jitter and then
reported as the derivative of the original model.

\begin{table}[ht]
\centering
\caption{Recommended route boundaries.}
\begin{tabular}{p{0.28\linewidth}p{0.29\linewidth}p{0.32\linewidth}}
\toprule
Situation & Admitted operation & Status \\
\midrule
Full-rank, well-separated stack & Fixed QR and block conditional QR & Value and
  analytical score \\
Full-rank but badly conditioned & Direct QR plus pivot/growth diagnostics; RRQR
  preflight & Score only while fixed branch remains valid \\
Known structural rank defect & Rectangular factor; fixed support chart & Value;
  score only with support-chart derivative certificate \\
Unknown or changing rank & Direct RRQR/SVD of residual stack & Value-only or
  fail closed \\
Singular innovation & Support likelihood and support-coordinate block QR & Value;
  score conditional on fixed support \\
Negative covariance weights & Ordered signed updates/downdates or hyperbolic QR &
  Separate audited branch; not covered by nonnegative block QR \\
\bottomrule
\end{tabular}
\end{table}

\section{Testing and evidence requirements}

The following tests are necessary to distinguish an actual square-root filter
from a covariance-form filter with a different name.

\subsection{Unit tests for algebraic primitives}

For random full-rank stacks, verify QR reconstruction, positive diagonal, solve
residuals, and the derivative identities in equations~\eqref{eq:qr-reconstruction}.
Use diagonal scales spanning many orders of magnitude and compare against an
independent high-precision or SVD diagnostic.  Test zero columns, duplicate
columns, and rank-one stacks as explicit rank events, not as ordinary passing
cases.

For block QR, verify all three identities in equation~\eqref{eq:block-covariance-identities}
and verify that $R_{xx}\transpose$ agrees with an independent Schur-complement
reference on benign positive-definite fixtures.  The production kernel must not
assemble that reference covariance.

For rectangular factors, test invariance under a fixed orthogonal rotation only
in the linear case.  In a nonlinear sigma-point transform, document and test
that changing the factor gauge can change the approximation; this is a feature
of sigma-point rules, not a numerical bug.

\subsection{Degeneracy and likelihood tests}

Construct exact rank-deficient process, state, and innovation fixtures.  Check
that the value route emits finite support likelihoods when observations lie on
the support, rejects off-support observations, and reports the correct retained
rank and pseudodeterminant.  Compare with a small independent generalized-
inverse calculation.  Test near-singular families on both sides of the rank
threshold and verify that score admission stops at the declared branch event.

\subsection{Filter-level parity tests}

On affine linear models, compare the direct-factor filter with a conventional
Kalman implementation for means, innovation likelihoods, gains, and covariance
reconstructions.  Compare the direct residual-stack route with an independent
SVD-of-stack value route, not only with an SVD-of-covariance route.  For regular
nonlinear fixtures, compare old and new sigma-point values and scores under a
predeclared tolerance and seed.  Long-run tests should monitor minimum QR pivot,
relative pivot, support residual, retained rank, factor reconstruction residual,
and solve residual.

\subsection{Derivative tests}

For a fixed-rank, fixed-pivot window, compare analytical scores with centered
finite differences of the same finite value program.  Also compare the
differentiated factor reconstruction with a covariance derivative assembled
only at the test boundary.  Repeat with parameter-dependent process and
observation loadings.  Deliberately cross a pivot or rank threshold and assert
that the score route fails closed rather than returning an apparently smooth
number.

\section{What is novel, what is not, and recommended decision}

The following classification prevents overclaiming.

\begin{description}
  \item[Source-faithful.] Sigma-point moment construction; QR residual-stack
    square-root propagation; triangular least-squares solves; Cholesky update or
    downdate for signed weights; and the numerical motivation for avoiding
    covariance square-rooting, all supported by
    \cite{vandermerwewan2001,arasaratnamhaykin2008,arasaratnamhaykin2009,
    kaminski1971,bierman1977}.
  \item[Compositional extension.] Rectangular factor coordinates combined with
    a direct residual-stack UKF and a support-aware degenerate innovation
    likelihood.  These follow directly from the factor and generalized-inverse
    mathematics, but the cited papers do not provide this exact end-to-end
    algorithm.
  \item[New implementation proposal.] The block-QR conditional factor in
    equation~\eqref{eq:conditional-factor}, used as the primary update so that the
    filtered factor is obtained by orthogonal elimination rather than
    covariance subtraction or sequential downdates.  It is a Schur-complement
    identity, not a claim that a particular paper published this exact API.
\end{description}

The recommendation is therefore:
\begin{enumerate}
  \item retain the current direct QR route as the differentiable default only
        for a declared fixed-rank, positive-pivot branch;
  \item replace its covariance-subtraction/downdate measurement update with the
        block conditional QR route where the nonnegative DZ5 weights apply;
  \item add a rectangular-factor and support-aware value route for structural
        singularity;
  \item use direct-stack RRQR/SVD as an independent rank/conditioning monitor
        and value-only fallback, never as an unqualified smooth score backend;
        and
  \item admit analytical scores only when rank, pivot pattern, sign convention,
        support chart, and likelihood measure remain fixed over the differentiated
        neighborhood.
\end{enumerate}

This is robust in the sense the square-root literature supports: it avoids Gram
matrix formation, preserves PSD factors by construction on the active branch,
uses stable orthogonal transformations and solves, and makes degeneracy visible.
It is not a promise that a single algorithm is smooth through arbitrary rank
changes.  No QR-versus-SVD choice can remove that mathematical fact.

\section{Repository model inventory and empirical coverage}
\label{sec:empirical-coverage}

The repository-wide coverage campaign was executed from the versioned artifact
root \texttt{docs/plans/artifacts/direct-factor-srukf-model-coverage-20260817/}.
The inventory contains 24 rows: every Common V2 model, every active and blocked
NeuTra registry cell, the explicit SSL--LSTM exclusion, the legacy actual-SV
route, and the generic MacroFinance adapter. Classification is a contract
classification, not a claim that every model is an SR--UKF problem. A
principal-root UKF, SGQF, multiplicative stochastic-volatility,
domain-clipped SIR, or covariance-form provider is not silently converted into
direct-factor evidence.

\begin{longtable}{p{0.29\linewidth}p{0.18\linewidth}p{0.45\linewidth}}
\caption{Generated inventory classification. The complete source anchors,
checksums, and dimensions are in \texttt{model\_inventory.json}.}\label{tab:inventory}\\
\toprule
Model or registry cell & Status & Contract boundary or reason \\
\midrule
\endfirsthead
\toprule
Model or registry cell & Status & Contract boundary or reason \\
\midrule
\endhead
\texttt{model\_a\_affine} & eligible\_score & Existing affine \texttt{TFFactorSRUKFModel} fixture. \\
\texttt{model\_b\_nonlinear\_accumulation} & eligible\_score & Existing nonlinear direct-factor fixture. \\
\texttt{model\_c\_nonlinear\_growth} & eligible\_score & Existing nonlinear direct-factor fixture. \\
\texttt{structural\_ar1\_quadratic\_h16} & eligible\_value\_only & Deterministic completion creates structural support; no score claim. \\
\texttt{lgssm\_2d\_h25\_rich} & eligible\_score & Common V2 affine factor adapter; exact SVD/Kalman authority. \\
\texttt{range\_bearing\_4d\_h20\_rich} & eligible\_score & Circular mean/residual adapter on a fixed branch. \\
\texttt{predator\_prey\_rk4} & eligible\_score & Common V2 physical-$r$ RK4 adapter; horizon $T=3$. \\
\texttt{PP-UKF}, \texttt{STR-UKF} & eligible\_score & One-time pre-trace factor adapters; temporal route is direct block-QR. \\
\texttt{LGSSM-EXACT} & eligible\_score & Full raw-18 likelihood and score tied to exact SVD authority. \\
\texttt{sv\_1d\_h18\_rich} & not\_applicable\_contract & Multiplicative observation $y=\exp(h/2)e$. \\
\texttt{spatial\_sir\_j3\_rk4} & not\_applicable\_contract & Domain-constrained epidemiological contract. \\
\texttt{PP-SGQF}, \texttt{SIR-SGQF} & not\_applicable\_contract & Different sparse-grid quadrature contract. \\
\texttt{SVX-ZC} & not\_applicable\_contract & Active frozen T10 Zhao--Cui TT target; GPU/XLA/HMC capable, but not direct-factor SR--UKF. \\
\texttt{macrofinance\_lgssm\_adapter} & not\_applicable\_contract & Generic external covariance-form provider. \\
\texttt{SVX-SGQF}, \texttt{KSC-UKF}, \texttt{PP-ZC}, \texttt{STR-ZC}, \texttt{SIR-ZC} & blocked & Current registry blockers preserved. \\
\texttt{actual\_sv\_independent\_panel} & historical\_only & Legacy manual factor route, not admitted block-QR contract. \\
\texttt{SIR-UKF}, \texttt{SSL-LSTM} & owner\_excluded & Explicit owner/user exclusions. \\
\bottomrule
\end{longtable}

Six direct-factor campaign rows were executable under the frozen contract
gate. The five score-bearing rows were run in eager and host-XLA modes, compared with the
historical principal-root implementation, and checked with direct QR pivot and
factor/derivative reconstruction diagnostics. PP-UKF and STR-UKF factor their
initial/noise covariance contracts once before tracing; no covariance is
factored inside their temporal recursions. The structural row was run by
the rectangular direct-stack SVD support route; it was deliberately not given
an analytical score.

\begin{table}[ht]
\centering
\caption{Generated direct block-QR score-bearing parity results. Errors compare
the new route with the historical principal-root route; nonlinear differences
depend on factor gauge and sigma-point orientation.}
\label{tab:score-parity}
\begin{tabular}{lrrrrrr}
\toprule
Model & $T$ & $|\Delta \ell|_{\max}$ & $|\Delta s|_{\max}$ & min pivot & factor residual & XLA \\
\midrule
\texttt{model\_a\_affine} & 5 & $0.0$ & $2.22\times10^{-16}$ & $1.27\times10^{-1}$ & $1.51\times10^{-16}$ & pass \\
\texttt{model\_b\_nonlinear\_accumulation} & 3 & $1.69\times10^{-4}$ & $8.42\times10^{-4}$ & $7.05\times10^{-2}$ & $1.15\times10^{-16}$ & pass \\
\texttt{model\_c\_nonlinear\_growth} & 2 & $3.00\times10^{-2}$ & $2.47\times10^{-2}$ & $2.80\times10^{-1}$ & $2.40\times10^{-15}$ & pass \\
\texttt{PP-UKF} & 20 & $2.22\times10^{-6}$ & $7.78\times10^{-6}$ & $1.44$ & $1.63\times10^{-15}$ & pass \\
\texttt{STR-UKF} & 100 & $1.66\times10^{-1}$ & $1.58$ & $2.21\times10^{-1}$ & $6.95\times10^{-16}$ & pass \\
\bottomrule
\end{tabular}
\end{table}

The affine row agrees to floating-point precision. The nonlinear rows are
finite and internally differentiated, but their historical-route deltas are
not a claim of exact equality: the old and new filters use different factor
gauges, and nonlinear sigma-point transforms are gauge-sensitive. The score
also agrees with centered finite differences of the same direct-factor value
program in the existing parity tests. For the two registry adapters, maximum
same-program score/finite-difference errors are $2.15\times10^{-9}$ (PP-UKF)
and $1.62\times10^{-8}$ (STR-UKF). The larger STR-UKF historical delta is
retained as factor-gauge-sensitive comparison evidence, not hidden as parity.

\begin{table}[ht]
\centering
\caption{Generated singular and ill-conditioned diagnostics.}
\label{tab:singular-robustness}
\begin{tabular}{lllrr}
\toprule
Case & Route & Result & Rank & Reconstruction residual \\
\midrule
Scales $1$ through $10^{-15}$ & direct QR & finite, positive-pivot branch & 2 & $0$ \\
Exact rank zero & direct-stack SVD & finite, value-only & 0 & $0$ \\
Exact rank one & direct-stack SVD & finite, value-only & 1 & $0$ \\
Repeated singular values & direct-stack SVD & finite, value-only & 2 & $0$ \\
Structural temporal fixture & rectangular support & on-support, XLA parity & 1 & $0$ \\
\bottomrule
\end{tabular}
\end{table}

The SVD rows are rank-discovery and support diagnostics only. Repeated
singular values, cutoff crossings, and support changes do not carry an
analytical score. The rectangular temporal fixture returned
$\ell=-1.4250004806894907$, minimum observation rank one, zero support
residual, and an eager/XLA absolute difference of
$8.88\times10^{-16}$.

\subsection{Empirical limitations and provenance}

These results establish a bounded numerical and model-contract campaign, not
exact nonlinear Bayesian inference, broad posterior correctness, HMC
readiness, or GPU production readiness. The six Common V2 rows and the full
NeuTra registry are classified. A superseding closure campaign certified the
four previously adapter-required rows under frozen full-rank/fixed-branch
contracts. No score is admitted for a singular or rank-changing branch. The
original campaign was generated by
\texttt{scripts/run\_direct\_factor\_srukf\_model\_coverage\_20260817.py}; its
manifest records the source revision, TensorFlow version, CPU/XLA diagnostic
lane, command lines, and artifact paths. Tables in this section were frozen
from \texttt{latex\_table\_payload.json}, with detailed values in the JSON
result files.

\subsection{Remaining-adapter closure campaign}
\label{sec:remaining-adapter-closure}

The four rows previously classified as \texttt{adapter\_required} were
certified by
\texttt{scripts/run\_direct\_factor\_srukf\_remaining\_adapter\_closure\_20260817.py}.
The final strict evidence root is
\texttt{docs/plans/artifacts/direct-factor-srukf-remaining-adapter-closure-20260817-r4/};
the earlier attempt roots are preserved as superseded harness provenance. The
superseding inventory contains nine \texttt{eligible\_score} rows, one
\texttt{eligible\_value\_only} row, and zero
\texttt{adapter\_required} rows. A later managed-GPU correction moved
\texttt{SVX-ZC} from the registry-blocked set to executable status for its
frozen T10 target. Within this direct-factor inventory it remains
\texttt{not\_applicable\_contract}, because a Zhao--Cui TT target is not an
SR--UKF contract.

\begin{table}[ht]
\centering
\caption{Remaining-adapter closure evidence. Authority deltas use an
independent linear-Gaussian SVD authority where applicable; FD is the same
finite direct-factor program's centered difference at $h=5\times10^{-6}$.}
\label{tab:remaining-adapter-closure}
\begin{tabular}{lrrrrrr}
\toprule
Model & $T$ & $|\Delta \ell|$ & $|\Delta s|$ & FD error & min pivot & XLA score \\
\midrule
\texttt{lgssm\_2d\_h25\_rich} & 25 & $0$ & $8.88\times10^{-16}$ & $1.03\times10^{-9}$ & $2.84\times10^{-1}$ & $1.78\times10^{-15}$ \\
\texttt{range\_bearing\_4d\_h20\_rich} & 20 & N/A & N/A & $5.68\times10^{-6}$ & $4.88\times10^{-2}$ & $4.83\times10^{-13}$ \\
\texttt{predator\_prey\_rk4} & 3 & N/A & N/A & $5.44\times10^{-9}$ & $1.45$ & $5.33\times10^{-14}$ \\
\texttt{LGSSM-EXACT} & 120 & $3.73\times10^{-14}$ & $7.46\times10^{-14}$ & $5.29\times10^{-9}$ & $8.06\times10^{-2}$ & $4.26\times10^{-14}$ \\
\bottomrule
\end{tabular}
\end{table}

The Common V2 LGSSM and LGSSM-EXACT rows agree with independent
linear-Gaussian authorities to floating-point precision. LGSSM-EXACT compares
likelihood first and adds the identical persisted Gaussian prior only
afterward. Predator-prey uses the Common V2 physical coordinate $r$ and its
three-observation fixture, not the separate PP-UKF probit fixture.
Range/bearing uses a circular bearing mean, wrapped sigma deviations, and
wrapped innovation. Its nominal minimum angular/resultant branch margin was
$0.988$. The centered-difference error decreased from
$2.27\times10^{-5}$ at $h=10^{-5}$ to $5.68\times10^{-6}$ at
$h=5\times10^{-6}$, the expected second-order convergence.

These closures do not change the singular/rank-changing policy. The direct
factor analytical score remains a fixed-positive-pivot and, where applicable,
fixed-angular-branch claim. Exact rank loss, pivot changes, and angular-cut
crossings remain value-only or score-blocked diagnostics.

\section{References}

\section{Remaining-gap campaign disposition}

The 17 August 2026 bounded campaign tested the outstanding repository cells
without conflating their target families. The KSC Gaussian-sum repair passed
the frozen T20 dense-value, score, mass-preservation, and permutation gates for
component caps $7,16,32,64,128,256$. Its GPU~3 XLA canary agreed with CPU to
$7.11\times10^{-15}$ in value and $2.22\times10^{-15}$ in score. This remains
a KSC surrogate result, not direct-factor SR--UKF, exact-SV, or HMC evidence.

The exact-SV SGQF ladder ran on GPU~3 at levels $10,12,16,20,24$, with level
$32$ as reference. No level passed: the smallest dense-prefix value error was
$3.38696\times10^{-3}$ per observation against the declared $10^{-3}$ gate.
This is a scientific negative result rather than an environment blocker.
PP--ZC and STR--ZC remain blocked until a source-anchored batch-native target
contract exists, and SIR--ZC remains blocked until an observed-data
parameter-score program is supplied.

The current frozen T10 SVX--ZC route passed GPU/CPU value and score parity at
$7.11\times10^{-15}$, same-program finite differences at
$4.62\times10^{-9}$, permutation invariance, status gates, XLA JIT, and the
memory-growth policy. Its target signature matches the prior terminal
sequential-HMC artifact and its current adapter signature matches the training
artifacts. The repository therefore restores
\texttt{xla\_hmc\_ready=true} and
\texttt{full\_chain\_xla\_diagnostic\_ready=true} for this frozen scope.
The score remains a manual same-program derivative, so
\texttt{runtime\_autodiff\_for\_hmc=false}; no exact-filter, broad-posterior,
superiority, or default-readiness claim follows.

These results reinforce the mathematical boundary of the direct-factor
construction.  A rectangular factor and support-density value can remain
well-defined at exact rank loss, but a single globally valid analytical score
cannot be asserted through rank changes, QR pivot/sign changes, repeated
singular values, or a zero circular resultant.  The implemented policy is
therefore value-only (or score-invalid) at those events, with fixed-rank,
fixed-pivot, fixed-branch charts required for score claims.  Full machine-
readable evidence is in
\path{docs/plans/artifacts/direct-factor-srukf-remaining-gaps-hypotheses-20260817/}.

\begin{thebibliography}{99}

\bibitem{arasaratnamhaykin2008}
I.~Arasaratnam and S.~Haykin, ``Square-root quadrature Kalman filtering,''
\emph{IEEE Transactions on Signal Processing}, vol.~56, no.~6, pp.~2589--2593,
2008.

\bibitem{arasaratnamhaykin2009}
I.~Arasaratnam and S.~Haykin, ``Cubature Kalman filters,''
\emph{IEEE Transactions on Automatic Control}, vol.~54, no.~6, pp.~1254--1269,
2009, doi:10.1109/TAC.2009.2019800.

\bibitem{benisraelgreville2003}
A.~Ben-Israel and T.~N.~E. Greville, \emph{Generalized Inverses: Theory and
Applications}, 2nd ed., Springer, 2003.

\bibitem{bierman1977}
G.~J. Bierman, \emph{Factorization Methods for Discrete Sequential Estimation},
Academic Press, 1977.

\bibitem{bogachev1998}
V.~I. Bogachev, \emph{Gaussian Measures}, American Mathematical Society, 1998.

\bibitem{chan1987}
T.~F. Chan, ``Rank revealing QR factorizations,'' \emph{Linear Algebra and its
Applications}, vols.~88--89, pp.~67--82, 1987.

\bibitem{golubvanloan2013}
G.~H. Golub and C.~F. Van Loan, \emph{Matrix Computations}, 4th ed., Johns
Hopkins University Press, 2013.

\bibitem{gu1996}
M.~Gu and S.~C. Eisenstat, ``Efficient algorithms for computing a strong
rank-revealing QR factorization,'' \emph{SIAM Journal on Scientific Computing},
vol.~17, no.~4, pp.~848--869, 1996.

\bibitem{ionescu2015}
C.~Ionescu, O.~Vantzos, and C.~Sminchisescu, ``Training deep networks with
structured layers by matrix backpropagation,'' arXiv:1509.07838, 2015.

\bibitem{julieruhlmann1997}
S.~J. Julier and J.~K. Uhlmann, ``A new extension of the Kalman filter to
nonlinear systems,'' in \emph{Proceedings of AeroSense}, 1997.

\bibitem{julieruhlmann2004}
S.~J. Julier and J.~K. Uhlmann, ``Unscented filtering and nonlinear estimation,''
\emph{Proceedings of the IEEE}, vol.~92, no.~3, pp.~401--422, 2004.

\bibitem{kochyang1998}
K.~R. Koch and Y.~Yang, ``Robust Kalman filter for rank deficient observation
models,'' \emph{Journal of Geodesy}, vol.~72, no.~7, pp.~436--441, 1998.

\bibitem{kaminski1971}
P.~G. Kaminski, A.~E. Bryson, and S.~Schmidt, ``Discrete square root filtering:
A survey of current techniques,'' \emph{IEEE Transactions on Automatic Control},
vol.~AC-16, no.~6, pp.~727--736, 1971.

\bibitem{magnusneudecker1999}
J.~R. Magnus and H.~Neudecker, \emph{Matrix Differential Calculus with
Applications in Statistics and Econometrics}, 2nd ed., Wiley, 1999.

\bibitem{maybeck1979}
P.~S. Maybeck, \emph{Stochastic Models, Estimation, and Control}, vol.~1,
Academic Press, 1979.

\bibitem{stewartsun1990}
G.~W. Stewart and J.-G. Sun, \emph{Matrix Perturbation Theory}, Academic Press,
1990.

\bibitem{vandermerwewan2001}
R.~van der Merwe and E.~A. Wan, ``The square-root unscented Kalman filter for
state and parameter estimation,'' in \emph{Proceedings of ICASSP}, vol.~6,
pp.~3461--3464, 2001.

\bibitem{vandermerwe2004}
R.~van der Merwe, \emph{Sigma-Point Kalman Filters for Probabilistic Inference
in Dynamic State-Space Models}, Ph.D. dissertation, OGI School of Science and
Engineering, Oregon Health and Science University, 2004.

\bibitem{yang2019robust}
C.~Yang, W.~Shi, and W.~Chen, ``Robust M-M unscented Kalman filtering for
GPS/IMU navigation,'' accepted manuscript, 2019, doi:
10.1007/s00190-018-01227-5.  The local source is a pre-publication version; the
journal, volume, and page metadata were not inferred from it.

\bibitem{tseng2017robust}
C.-H. Tseng, S.-F. Lin, and D.-J. Jwo, ``Robust Huber-based cubature Kalman
filter for GPS navigation processing,'' \emph{Journal of Navigation}, vol.~70,
pp.~527--546, 2017, doi:10.1017/S0373463316000692.

\end{thebibliography}

\end{document}
